%%%%%%%%%%  scan idx 440  |  folio 427  %%%%%%%%%%
\SGAPrintPage{427}%
this kernel is contained in that of the composite homomorphism:
\[
\mathbb Z[T_1,X_1,\ldots,X_d]\longrightarrow
\bigoplus_{n\ge0}J^n\longrightarrow A,
\]
the last homomorphism being defined in degree \(n\) by the inclusion of \(J^n\) in \(A\). Now the kernel of the composite homomorphism is the principal ideal generated by \((X_1-T_1)\), which, for reasons of integrality, contains no nonzero homogeneous elements. Thus the first homomorphism is injective, completing the proof in this case.

Return to the general case where \(A\) is arbitrary; we may still suppose \(\underline{E}\) free over \(X\). Attach the index \(0\) to the special case just treated; then, by part i) of the proposition, the situation under consideration is obtained by the base change \(A_0\to A\) from the situation studied above, the homomorphism \(A_0\to A\) being defined by taking a basis of \(E\), its image \(f_1,\ldots,f_d\) in \(J\), and sending \(T_i\) to \(f_i\). Put \(B_0=\left(\bigoplus_{n\ge0}J_0^n\right)_{(T_1^{(1)})}\), \(B=\left(\bigoplus_{n\ge0}J^n\right)_{(f_1^{(1)})}\); then \(B=B_0\otimes_{A_0}A\); the Koszul complex defined by \(v:\underline{F}\to\underline{O}_P\) is deduced by base change from the Koszul complex defined by \(v_0\); since the latter is a projective resolution of \(B_0\) (working on a suitable open subset), the homology groups of \(K_\bullet(v)\) on this open subset are the \(\Tor_i^{A_0}(A,B_0)\). These are the degree \(0\) components of the \(\Tor_i^{A_0}\left(A,\left(\bigoplus_{n\ge0}J_0^n\right)_{T_1^{(1)}}\right)=\left(\bigoplus_{n\ge0}\Tor_i^{A_0}(A,J_0^n)\right)_{T_1^{(1)}}\); the \(\Tor_i^{A_0}(A,J_0^n)\) vanish for \(i\ge1\) by the relations 1.2, so \(K_\bullet(v)\) is a resolution of \(B\), which shows that \(j\) is regular.

iii) Since the property is local on \(X\), we may again suppose \(X\) affine and \(\underline{E}\) free of finite type. Resume the notation of ii).

First consider the standard case of ii). The ring \(A_0\) is then noetherian and of finite cohomological dimension, and the algebra \(\mathbf{S}(E_0)\)

%%%%%%%%%%  scan idx 441  |  folio 428  %%%%%%%%%%
\SGAPrintPage{428}%
likewise. We can then construct step by step a partial graded resolution of \(\bigoplus_{n\ge0}J_0^n\) by graded free \(\mathbf{S}(E_0)\)-modules of finite type, of length greater than the cohomological dimension of \(A_0\); the last kernel is then a graded \(\mathbf{S}(E_0)\)-module, projective of finite type as an ungraded \(\mathbf{S}(E_0)\)-module, hence locally free of finite type over \(\Spec(A_0)\) as a graded \(\mathbf{S}(\underline{E}_0)\)-module by VI 1.5.1. We obtain the announced finite resolution by replacing the last module of the partial resolution by the corresponding kernel.

Tensoring the resolution thus constructed with \(A\), we obtain a graded complex of graded \(\mathbf{S}(E)\)-modules locally free of finite type, and must show that it is a resolution of \(\bigoplus_{n\ge0}J^n\). The \(i\)-th homology Module of the resulting complex is a graded \(\mathbf{S}(E)\)-Module whose \(n\)-th component is \(\Tor_i^{A_0}(A,J_0^n)\), which vanishes for \(i\ge1\) by (1.2.4). Consequently, by (1.2.3), we have constructed a resolution of the required type.

\sgaStableTarget{sga6.e07.prop.1-9}\textbf{Proposition 1.9.} \uline{Let} \(i:Y\to X\) \uline{be a regular closed immersion}, \(X'\) \uline{the} \(X\)-\uline{scheme obtained by blowing up} \(Y\), \uline{and} \(f:X'\to X\) \uline{the structural morphism}. \uline{Then} \(i\) \uline{and} \(f\) \uline{are perfect morphisms}. (III 4.1.).

Since the property is local on \(X\), we may suppose given an \(\underline{O}_X\)-Module locally free of finite type \(\underline{E}\), and a regular homomorphism \(u:\underline{E}\to\underline{O}_X\) whose image is the ideal \(\underline{J}\) of \(Y\). The Koszul complex \(\underline{K}_\bullet(u)\) is then a finite resolution of \(\underline{O}_Y\) by \(\underline{O}_X\)-Modules locally free of finite type; consequently, \(\underline{O}_Y\) is a perfect \(\underline{O}_X\)-Module. It follows that every regular immersion is a perfect morphism. By 1.8, \(f\) can be written as the composite of a regular immersion \(X'\to\bfP(\underline{E})\) and the structural morphism \(\bfP(\underline{E})\to X\), which is smooth and hence perfect; it is therefore

%%%%%%%%%%  scan idx 442  |  folio 429  %%%%%%%%%%
\SGAPrintPage{429}%
the composite of two perfect morphisms, and consequently a perfect morphism.

\sgaStableTarget{sga6.e07.prop.1-10}\textbf{Proposition 1.10.} \uline{Let} \(X\), \(Y\) \uline{be two schemes smooth over} \(S\). \uline{Then every} \(S\)-\uline{immersion of} \(X\) \uline{into} \(Y\) \uline{is regular}. \uline{In particular}, \uline{every section of a smooth morphism is a regular immersion}.

The proposition follows immediately from EGA IV 17.12.1, since a regular immersion in the sense of EGA IV is regular in the sense of 1.4.

\sgaStableTarget{sga6.e07.sec.2}\section*{2. Computations on regular immersions}

\sgaStableTarget{sga6.e07.para.2-1}\subsection*{2.1.}
Let \(A\) be a ring, and \(B\) and \(C\) two \(A\)-algebras. We propose to endow \(\bigoplus_{i\ge0}\Tor_i^A(B,C)\) with a graded \(A\)-algebra structure, and hence to define homomorphisms:
\begin{equation}\tag{2.1.1}
 \Tor_i^A(B,C)\otimes_A\Tor_j^A(B,C)\longrightarrow \Tor_{i+j}^A(B,C).
\end{equation}

Let \(P_\bullet\) be a resolution of \(B\) by projective \(A\)-modules; the \(\Tor_i^A(B,C)\) are the homology groups of the complex \(P_\bullet\otimes_A C\). We know how to define canonical homomorphisms ([1], IV 6):
\begin{equation}\tag{2.1.2}
 H_i(P_\bullet\otimes_A C)\otimes_A H_j(P_\bullet\otimes_A C)
 \longrightarrow H_{i+j}\bigl((P_\bullet\otimes_A C)\otimes_A(P_\bullet\otimes_A C)\bigr).
\end{equation}

Furthermore, \(P_\bullet\otimes_A P_\bullet\) is a complex of projective \(A\)-modules, in degrees \(\geq0\), and \(H_0(P_\bullet\otimes_A P_\bullet)=B\otimes_A B\); moreover, \(P_\bullet\) is a resolution of \(B\). Consequently, the augmentation morphism \(B\otimes_A B\to B\) extends to a morphism of complexes: \(P_\bullet\otimes_A P_\bullet\to P_\bullet\) and therefore defines, using the homomorphism \(C\otimes_A C\to C\) deduced from the algebra structure of \(C\), homomorphisms:
\begin{equation}\tag{2.1.3}
 H_{i+j}\bigl((P_\bullet\otimes_A C)\otimes_A(P_\bullet\otimes_A C)\bigr)
 \longrightarrow H_{i+j}(P_\bullet\otimes_A C).
\end{equation}

%%%%%%%%%%  scan idx 443  |  folio 430  %%%%%%%%%%
\SGAPrintPage{430}%
Composing (2.1.2) and (2.1.3), we obtain the homomorphisms (2.1.1).

The algebra structure defined by the homomorphisms (2.1.1) is associative (cf. [1] XI 2). Furthermore, the \(A\)-module \(H_i(P_\bullet\otimes_A C)\) has a standard \(B\otimes_A C\)-module structure, with \(C\) acting in the evident way, and \(B\) by functoriality; likewise, the \(H_j\bigl((P_\bullet\otimes_A C)\otimes_A(P_\bullet\otimes_A C)\bigr)\) carry two \(B\otimes_A C\)-module structures. The homomorphisms (2.1.2) and (2.1.3) are clearly compatible with these two structures: for (2.1.2), by functoriality (working in the category of \(A\)-modules), and for (2.1.3), by the classical result on morphisms of an acyclic complex. We have therefore endowed \(\bigoplus_{i\ge0}\Tor_i^A(B,C)\) with an associative \(B\otimes_A C\)-algebra structure.

\subsection*{2.2.}
Suppose that \(B\) admits a resolution by projective \(A\)-modules \(P_\bullet\) such that \(P_\bullet=\bigoplus_{i\ge0}P_i\) is a graded \(A\)-algebra, with the differential satisfying:
\[
\begin{array}{ll}
\text{a)}& P_0\longrightarrow B\text{ is a ring homomorphism~;}\SGAAuthorFootMark{1}\\[0.4ex]
\text{b)}& \forall x\in P_i,\ \forall y\in P_j,\quad d(xy)=d(x)y+(-1)^i x\,d(y).
\end{array}
\]
\SGAAuthorFootText{1}{One can show that such a resolution always exists.}
Then the ``product'' maps define commutative diagrams for \(i,j\ge0\):
\sgaStableTarget{sga6.e07.eq.2-2-1}\begin{equation}\notag
\begin{tikzcd}[column sep=large,row sep=large]
P_i\otimes_A P_j \arrow[r] \arrow[d] & (P_{i-1}\otimes_A P_j)\oplus(P_i\otimes_A P_{j-1}) \arrow[d]\\
P_{i+j} \arrow[r] & P_{i+j-1},
\end{tikzcd}
\end{equation}
and consequently a morphism of complexes \(P_\bullet\otimes_A P_\bullet\to P_\bullet\) extending the augmentation homomorphism: \(B\otimes_A B\to B\). This shows that the algebra structure of \(\bigoplus_{i\ge0}\Tor_i^A(B,C)\) can be computed from that of \(P_\bullet\). In particular, if \(P_\bullet\) is an alternating \(A\)-algebra (i.e. such that for every

%%%%%%%%%%  scan idx 444  |  folio 431  %%%%%%%%%%
\SGAPrintPage{431}%
\(x\in P_\bullet\), we have \(x^2=0\)), then the same is true of \(\bigoplus_{i\ge0}\Tor_i^A(B,C)\).

This last remark applies when \(B\) is the quotient of \(A\) by a regular ideal, since we can then take for \(P_\bullet\) the Koszul complex associated with a regular homomorphism defining the ideal of \(B\).

\subsection*{2.3.}
The \(B\otimes_A C\)-algebra structure defined in 2.1 on \(\bigoplus_{i\ge0}\Tor_i^A(B,C)\) is functorial in \(B\) and \(C\). We can therefore globalize the preceding constructions: let \(X\) be a scheme, \(X'\), \(X''\) two \(X\)-schemes; we can then define the morphisms (2.1.1) locally on \(X'\times_X X''\); they glue by functoriality and define on \(\bigoplus_{i\ge0}\underline{\Tor}_i^{\underline{O}_X}(\underline{O}_{X'},\underline{O}_{X''})\) an associative graded \(\underline{O}_{X'\times_X X''}\)-Algebra structure.

\sgaStableTarget{sga6.e07.lem.2-4}\textbf{Lemma 2.4.} \uline{Let} \(X',X''\) \uline{be two} \(X\)-\uline{schemes}, \uline{and} \(\underline{M}\) \uline{a flat} \(\underline{O}_{X'}\)-\uline{Module}. \uline{Then for every} \(i\ge0\), \uline{we have a canonical isomorphism}:
\sgaStableTarget{sga6.e07.eq.2-4-1}\begin{equation}\notag
\underline{\Tor}_i^{\underline{O}_X}(\underline{M},\underline{O}_{X''})\simeq
\underline{M}\otimes_{\underline{O}_{X'}}\underline{\Tor}_i^{\underline{O}_X}(\underline{O}_{X'},\underline{O}_{X''}).
\end{equation}

It suffices to define a functorial isomorphism in the affine case; we can then define it globally by gluing in the general case.

Let \(X=\Spec(A)\), \(X'=\Spec(A')\), \(X''=\Spec(A'')\); \(\underline{M}\) is defined by a flat \(A'\)-module \(M\). Let \(P_\bullet\) be a projective resolution of \(A''\) over \(A\); the \(\Tor_i^A(A',A'')\) are the homology groups of \(A'\otimes_A P_\bullet\). Since \(M\) is flat over \(A'\), we have a canonical isomorphism: \(M\otimes_{A'}\Tor_i^A(A',A'')\xrightarrow{\sim}H_i\bigl(M\otimes_{A'}(A'\otimes_A P_\bullet)\bigr)\). Furthermore, the canonical isomorphism: \(M\otimes_{A'}(A'\otimes_A P_\bullet)\xrightarrow{\sim}M\otimes_A P_\bullet\) gives a canonical isomorphism: \(H_i\bigl(M\otimes_{A'}(A'\otimes_A P_\bullet)\bigr)\xrightarrow{\sim}\Tor_i^A(M,A'')\), whence the result.

\sgaStableTarget{sga6.e07.prop.2-5}\textbf{Proposition 2.5.} \uline{Let} \(X\) \uline{be a scheme}, \(Y\) \uline{and} \(Z\) \uline{two closed subschemes of} \(X\), \(T=Y\times_X Z\). \uline{Suppose that for every} \(x\in T\), \uline{there exists an open neighborhood}
